% ⚠️ 本章多处原文未写完（2026-10-20 逐行核对 md 与 tex 后确认，**两份在同样位置都残缺**）：
%   ① 定义7.2 之后的「对 / 不含 / 通解」三行是残句（特解 / 通解 / 积分曲线的定义各只写了一半）；
%   ② 一阶齐次线性方程「分离变量得」之后为空；
%   ③ 定理7.3、7.4、7.6、7.7、7.8、7.9、7.10 的陈述只写了开头半句；
%   ④ 可降阶高阶方程的三种类型（y''=f(x,y')、y''=f(y,y')）只有标题；
%   ⑤ 特征根为重根的情形只有「当有重特征根时，」；
%   ⑥ 欧拉方程「则由上述结果得到」之后为空。
%   补齐须对照教材，**不要凭猜测填写**。

% 子文件：7.微分方程.tex
% 主文档使用 \input{7.微分方程} 引入，无 document 环境

\section{微分方程的基本概念}

定义7.1\quad 含有未知函数的导数或微分的，联系着自变量、未知函数及其导数或微分的方程叫做微分方程。

微分方程中可以不明显出现自变量和未知函数，但必须含有未知函数的导数或微分。

在微分方程中出现的未知函数的导数的最高阶数，称为微分方程的阶。

$n$阶微分方程的一般形式为
\[
F(x,y,y',\cdots,y^{(n)})=0 \tag{1}
\]
这是联系着$x,y,y',\cdots,y^{(n)}$的关系式，式中$y^{(n)}$必须出现。

定义7.2\quad 凡满足微分方程的函数（即把它和它的导数代入微分方程后，能使方程变为恒等式）都叫做该微分方程的解。

对

不含

微分方程的特解的几何图形叫做积分曲线。

通解

定义7.3\quad 当自变量取某值时，要求未知函数及其导数取定值的条件叫做初值条件（初始条件）。

$n$阶微分方程的通解中含有$n$个任意常数，要想得到确定的解，必须且只需有$n$个定解条件，所以$n$阶微分方程的初值条件应为
\[
y(x_0)=y_0,\ y'(x_0)=y_0',\ \cdots,\ y^{(n-1)}(x_0)=y_0^{(n-1)}, \tag{2}
\]
其中$y_0,y_0',\cdots,y_0^{(n-1)}$是$n$个数值。

求微分方程满足初值条件的解的问题，叫做初值问题或柯西问题。

\section{一阶微分方程}
\subsection{可分离变量的方程}

如果一阶微分方程可以写成$g(y)dy=f(x)dx$的形式，其中等式的每一边仅是一个变量的函数与该变量的微分，则说这个方程是可分离变量的方程。

如
\[
\frac{dy}{dx}=f(x)h(y)
\]
和
\[
M_1(x)M_2(y)dx+N_1(x)N_2(y)dy=0
\]
都是可分离变量的方程。

\noindent 分离变量法

可分离变量的方程求通解的步骤是：
\begin{enumerate}
    \item 分离变量。
    \item 将上式两边积分
\[
\int g(y)dy=\int f(x)dx+C,
\]
其中$C$为任意常数。由上式确定的函数$y=y(x,C)$就是方程的通解。
\end{enumerate}

\subsection{一阶线性微分方程}

形如
\[
\frac{dy}{dx}+P(x)y=Q(x)
\]
的方程称为一阶线性微分方程。它是未知函数及其导数的一次方程，其中$P(x),Q(x)$为某区间上的已知函数。当$Q(x)\equiv0$时，方程
\[
\frac{dy}{dx}+P(x)y=0
\]
称为一阶齐次线性方程，相应地把$Q(x)\not\equiv0$的方程称为一阶非齐次线性方程。

一阶齐次线性方程是可分离变量的方程，分离变量得

积分并化简得方程的通解为

其中$C$为任意常数。

齐次线性方程通解作常数变易型，代入方程，得

从而$C(x)$满足方程

积分得

于是一阶非齐次线性方程的通解公式为
\[
y=e^{-\int P(x)dx}\Bigl(C+\int Q(x)e^{\int P(x)dx}dx\Bigr).
\]
把齐次线性方程的通解。

\subsection{变量代换}
\begin{enumerate}
    \item 如果一阶微分方程可以写成
\[
\frac{dy}{dx}=g\left(\frac{y}{x}\right)
\]
的形式，则称之为齐次方程。

做变换 $u=\dfrac{y}{x}$，即 $y=ux$，这是可分离变量的方程。

    \item 形如
\[
\frac{dy}{dx}+P(x)y^{1-n}=Q(x)
\]
的方程叫做伯努利方程。

事实上，即
\[
\frac{1}{1-n}\frac{d(y^{1-n})}{dx}+P(x)y^{1-n}=Q(x).
\]
可见只要作变换，令$z=y^{1-n}$，方程就可化为$z$的一阶线性方程
\[
\frac{dz}{dx}+(1-n)P(x)z=(1-n)Q(x).
\]
求出
\end{enumerate}

\section{几种可积的高阶微分方程}

$y^{(n)}=f(x)$型方程

只需积分$n$次。

$y''=f(x,y')$型方程

$y''=f(y,y')$型方程

\section{线性微分方程及其通解结构}

一般地，形如
\[
y^{(n)}+a_1(x)y^{(n-1)}+\cdots+a_n(x)y=f(x)
\]
的方程称为$n$阶线性微分方程。当$f(x)\equiv0$时，
\[
y^{(n)}+a_1(x)y^{(n-1)}+\cdots+a_n(x)y=0
\]
称为齐次线性微分方程。

\noindent 线性微分方程通解的结构

定理7.1（叠加原理）\quad 设函数$y_1(x)$与$y_2(x)$是齐次方程的两个解，则$y=C_1y_1(x)+C_2y_2(x)$也是方程的解，其中$C_1,C_2$是任意常数。

设$y_1(x),y_2(x),\cdots,y_n(x)$为定义在区间$I$上的函数，如果存在$n$个不全为零的常数，使得

则称这$n$个函数在区间$I$上线性相关，否则称线性无关。

现考虑区间$I$上的两个线性相关函数。使得

在区间$I$上恒成立，又得

这两个等式组成，所以应有。

$W(x)$称为函数组的朗斯基行列式。如果是齐次方程的解，容易证明解组线性无关的充要条件是朗斯基行列式$W(x)\neq0$。

定理7.2\quad 如果$y_1(x)$与$y_2(x)$是二阶齐次方程的两个线性无关的特解，则$y=C_1y_1(x)+C_2y_2(x)$是方程的通解，其中$C_1,C_2$是任意常数。

定理7.3\quad 设$y^{*}(x)$是非齐次线性微分方程
\[
y''+p(x)y'+q(x)y=f(x)
\]
的一个特解。

定理7.4\quad 设$y_1^{*}(x)$与$y_2^{*}(x)$依次是方程

的解，则$y=y_1^{*}(x)+y_2^{*}(x)$是方程

的解。

这一定理通常称为非齐次线性微分方程解的叠加原理。

定理7.5(齐次线性微分方程通解结构)\quad 设$y_1(x),y_2(x),\cdots,y_n(x)$是$n$阶齐次线性微分方程的$n$个线性无关的解，则方程的通解为
\[
y(x)=C_1y_1(x)+C_2y_2(x)+\cdots+C_ny_n(x),
\]
其中$C_1,C_2,\cdots,C_n$为任意常数。

定理7.6(非齐次线性微分方程通解结构)\quad 设$y^{*}(x)$是非齐次方程的一个特解。

\section{常系数线性微分方程}
\subsection{常系数齐次线性微分方程}

当$a_i\ (i=1,2,\cdots,n)$均为常数时，称
\[
y^{(n)}+a_1y^{(n-1)}+\cdots+a_ny=0
\]
为常系数齐次线性微分方程。

称
\[
r^n+a_1r^{n-1}+\cdots+a_n=0
\]
为方程的特征方程，并称其根为特征根。

\begin{enumerate}
    \item 当特征根为互不相同的实根时，就是常系数齐次线性微分方程的一个解组。

当出现共轭复根时，由欧拉公式
\[
e^{\alpha+i\beta}=e^\alpha(\cos\beta+i\sin\beta),\quad e^{\alpha-i\beta}=e^\alpha(\cos\beta-i\sin\beta)
\]
知，解

根据叠加原理知

是对应的两个实值解。

    \item 当有重特征根时，
\end{enumerate}

\subsection{常系数非齐次线性微分方程}

当$a_i\ (i=1,2,\cdots,n)$均为常数时，称
\[
y^{(n)}+a_1y^{(n-1)}+\cdots+a_ny=f(x)
\]
为常系数非齐次线性微分方程。

\noindent 待定系数法

方程的右端函数$f(x)$是多项式、正弦函数、余弦函数、指数函数。

设方程的右端函数
\[
f(x)=e^{\alpha x}\bigl[P(x)\cos \beta x+Q(x)\sin \beta x\bigr],
\]
其中$P(x),Q(x)$是多项式，$m$是它们的次数的最大值。

\subsection{欧拉方程}

形如
\[
x^ny^{(n)}+a_1x^{n-1}y^{(n-1)}+\cdots+a_{n-1}xy'+a_ny=f(x)
\]
的方程称为欧拉方程，其中$a_i\ (i=1,2,\cdots,n)$均为常数。

如果采用算子$D$表示关于$t$的导数运算$\dfrac{d}{dt}$，则由上述结果得到。

\section{线性微分方程组}
\subsection{线性微分方程组}

对于$n$阶线性微分方程
\[
y^{(n)}+a_1(x)y^{(n-1)}+\cdots+a_n(x)y=f(x),
\]
若令$y'=y_1,\ y''=y_2,\ \cdots,\ y^{(n-1)}=y_{n-1}$，则化为等价的方程组。

形如
\[
\begin{cases}
y_1'=a_{11}(x)y_1+\cdots+a_{1n}(x)y_n+f_1(x),\\
\quad\vdots\\
y_n'=a_{n1}(x)y_1+\cdots+a_{nn}(x)y_n+f_n(x)
\end{cases}
\]
的微分方程组称为（标准的）一阶线性微分方程组。通解是含有$n$个独立任意常数。式的初值条件为

其中$y_i^{[0]}\ (i=1,2,\cdots,n)$为给定的常数。

进一步，为了表述简洁，我们用向量和矩阵来表示，令

并规定向量函数和矩阵函数的连续、导数和积分是对它们的每个元素而言的。则可简记为
\[
\boldsymbol{y}'=\boldsymbol{A}(x)\boldsymbol{y}+\boldsymbol{f}(x).
\]

若$\boldsymbol{f}(x)\not\equiv \boldsymbol{0}$，称为非齐次线性微分方程组，$\boldsymbol{f}(x)$称为非齐次项；当$\boldsymbol{f}(x)\equiv \boldsymbol{0}$时，
\[
\boldsymbol{y}'=\boldsymbol{A}(x)\boldsymbol{y}
\]
称为齐次线性微分方程组。

\subsection{线性微分方程组通解结构}

定理7.7（存在唯一性定理）\quad 若$\boldsymbol{A}(x),\boldsymbol{f}(x)$在区间$I$上连续，则对任一。

定理7.8（叠加定理）\quad 设$\boldsymbol{y}_1,\boldsymbol{y}_2$依次是方程组

的解，则$a\boldsymbol{y}_1+b\boldsymbol{y}_2$是方程组

的解，其中$a,b$为常数。

定理7.9（齐次线性微分方程组通解结构）\quad 若齐次方程组有$n$个线性无关的解。

定理7.10（非齐次线性微分方程组通解结构）\quad 设$\boldsymbol{y}^*(x)$是非齐次线性微分方程组的一个特解。

\subsection{常系数齐次线性微分方程组}

当$\boldsymbol{A}$为$n\times n$阶常数矩阵时，称
\[
\boldsymbol{y}'=\boldsymbol{A}\boldsymbol{y}
\]
为常系数齐次线性微分方程组。设想有形如
\[
\boldsymbol{y}=\boldsymbol{v}e^{\lambda x}
\]
的解。